\subsection{双代数的本原元素}

设 E 为双代数（《代数》，第 III 章，第 11 节，第 4 号），c 为其余积，ε 为其余单位，1 为其单位元素。由于 ε(1) = 1 且 c(1) = 1 ⊗ 1，前一号的结果可在 u = 1 下应用。E 的 1-本原元素（第 1 号，定义 1）简称为本原元素（参看~《代数》，第 III 章，第 11 节，第 8 号），即 E 中满足下式的元素 x：

\[
c (x) = x \otimes 1 + 1 \otimes x.\tag{5}
\]

我们用 \(\pi\)、\(\eta\)、P(E)、\(c^{+}\) 简记 \(\pi_{1}\)、\(\eta_{1}\)、P \(_{1}\) (E)、\(c_{1}^{+}\)。

命题 4. 本原元素的集合 \(\mathrm{P(E)}\)（\(\mathbf{E}\) 的）是 \(\mathbf{E}\) 的李子代数。

若 \(x, y\) 属于 \(\mathrm{P(E)}\)，则

\[
\begin{array}{r l} c (x y) & = c (x) c (y) = (x \otimes 1 + 1 \otimes x) (y \otimes 1 + 1 \otimes y) \\ & = x y \otimes 1 + 1 \otimes x y + x \otimes y + y \otimes x, \end{array}
\]

由此

\[
c ([ x, y ]) = [ x, y ] \otimes 1 + 1 \otimes [ x, y ].
\]

命题 5. 设 \(f: \mathrm{E} \to \mathrm{E}'\) 为双代数态射。若 \(x\) 是 \(\mathrm{E}\) 的本原元素，则 \(f(x)\) 是 \(\mathrm{E}'\) 的本原元素，且 \(f\) 在 \(\mathrm{P(E)}\) 上的限制是李代数同态 \(\mathrm{P}(f): \mathrm{P(E)} \to \mathrm{P(E')}\)。

设 \(c\)（相应地，\(c'\)）为 E（相应地，E'）的余积。由于 \(f\) 是余代数态射，\(c' \circ f = (f \otimes f) \circ c\)，由此

\[
\begin{array}{r l} c ^ {\prime} (f (x)) & = (f \otimes f) (c (x)) = (f \otimes f) (x \otimes 1 + 1 \otimes x) \\ & = f (x) \otimes 1 + 1 \otimes f (x), \end{array}
\]

其中 \(x\) 是本原的。因而 \(f\) 把 \(\mathrm{P(E)}\) 映入 \(\mathrm{P(E')}\)，且 \(f([x,y]) = [f(x),f(y)]\)，因为 \(f\) 是代数同态。

注. (1) 设 \(p\) 为素数，使得 \(p\) . \(1 = 0\) 在 \(\mathbf{K}\) 中成立。二项式公式以及同余关系 \(\binom{p}{i} \equiv 0 \pmod{p}\)（对 \(1 \leqslant i \leqslant p - 1\)）蕴含

P(E) 在映射 \(x \mapsto x^p\) 下稳定。

\begin{enumerate}
\def\labelenumi{(\arabic{enumi})}
\setcounter{enumi}{1}
\tightlist
\item
  按定义，图示
\end{enumerate}

\[
0 \longrightarrow \mathrm{P(E)} \longrightarrow \mathrm {E^ {+}} \stackrel {c ^ {+}} {\longrightarrow} \mathrm {E^ {+}} \otimes \mathrm {E^ {+}}
\]

是正合列。若 \(\mathbf{K}'\) 是交换环且 \(\rho: \mathbf{K} \to \mathbf{K}'\) 是环同态，则 \(\rho^*(\mathrm{E}) = \mathrm{E} \otimes_{\mathbb{K}} \mathbf{K}'\) 是 \(\mathbf{K}'\) -双代数，且含入 \(\mathrm{P}(\mathrm{E}) \to \mathrm{E}\) 定义了李 \(\mathbf{K}'\) -代数的同态

\[
\alpha \colon \mathrm{P} (\mathrm{E}) \otimes_ {\mathrm{K}} \mathrm{K} ^ {\prime} \rightarrow \mathrm{P} (\mathrm{E} \otimes_ {\mathrm{K}} \mathrm{K} ^ {\prime}).
\]

若 \(\mathbf{K}'\) 在 \(\mathbf{K}\) 上是平坦的（《交换代数》，第 I 章，第 2 节，第 3 号，定义 2），则由同处可知图示

\[
0 \longrightarrow \mathrm{P} (\mathrm{E}) \otimes_ {\mathrm{K}} \mathrm{K} ^ {\prime} \longrightarrow \mathrm{E} ^ {+} \otimes_ {\mathrm{K}} \mathrm{K} ^ {\prime} \xrightarrow {c ^ {+} \otimes_ {\mathrm{K}} \mathrm{Id} _ {\mathrm{K} ^ {\prime}}} \left(\mathrm{E} ^ {+} \otimes_ {\mathrm{K}} \mathrm{K} ^ {\prime}\right) \otimes_ {\mathrm{K} ^ {\prime}} \left(\mathrm{E} ^ {+} \otimes_ {\mathrm{K}} \mathrm{K} ^ {\prime}\right)
\]

是正合列，由此推出 \(\alpha\) 是同构。

